% Section: truth-in-model

\documentclass[../../../include/open-logic-section]{subfiles}

\begin{document}

\olfileid{nml}{syn}{tru}

\olsection{মডেলে সত্যতা}

কখনো আমরা জানতে চাই, কোনো নির্দিষ্ট মডেলের প্রত্যেক জগতে কোন !!{formula}s সত্য।
এর জন্য একটি সংকেত নির্ধারণ করি।

\begin{defn}
  !!^a{formula}~$!A$ মডেল $M = \tuple{W, R,
    V}$-এ \emph{সত্য}, সংকেতে $\mSat{M}{!A}$, তখনই এবং কেবল তখনই, যখন
  প্রত্যেক $w \in W$-এর জন্য $\mSat{M}{!A}[w]$।
\end{defn}

\begin{prop}\ollabel{prop:truthfacts}
  \begin{enumerate}
  \item $\mSat{M}{!A}$ হলে $\mSat/{M}{\lnot !A}$; কিন্তু এর বিপরীতটি
    \emph{সত্য নয়}।
  \item $\mSat{M}{!A \lif !B}$ হলে $\mSat{M}{!A}$ থেকে
    $\mSat{M}{!B}$ পাওয়া যায়; কিন্তু এর বিপরীতটি \emph{সত্য নয়}।
  \end{enumerate}
\end{prop}

\begin{proof}
  \begin{enumerate}
  \item $\mSat{M}{!A}$ হলে $!A$, $W$-এর সব জগতে সত্য। যেহেতু
    $W \neq \emptyset$, $\mSat{M}{\lnot!A}$ হতে পারে না; তা হলে কোনো
    জগতে $!A$-কে একই সঙ্গে সত্য ও অসত্য হতে হতো।

    অন্যদিকে, $\mSat/{M}{\lnot !A}$ হলে $!A$ কোনো $w \in W$-তে সত্য।
    তাই বলে \emph{প্রত্যেক}~$w \in W$-এ $\mSat{M}{!A}[w]$ হয় না।
    যেমন, \olref[rel]{fig:simple}-এর মডেলে $\mSat/{M}{\lnot p}$,
    আবার $\mSat/{M}{p}$-ও।
  \item ধরা যাক, $\mSat{M}{!A \lif !B}$ এবং $\mSat{M}{!A}$।
    $\mSat{M}{!B}$ দেখাতে যেকোনো $w \in W$ নিই। তখন
    $\mSat{M}{!A \lif !B}[w]$ এবং $\mSat{M}{!A}[w]$; তাই
    $\mSat{M}{!B}[w]$। $w$ নির্বিচারে নেওয়া হয়েছে বলে
    $\mSat{M}{!B}$।

    বিপরীতটি যে সত্য নয়, তা দেখাতে এমন একটি মডেল $\mModel{M}$ চাই যাতে
    $\mSat{M}{!A}$ হলে $\mSat{M}{!B}$, অথচ
    $\mSat/{M}{!A \lif !B}$। আবার \olref[rel]{fig:simple}-এর মডেলটি
    দেখো: $\mSat/{M}{p}$, তাই (শূন্যতাবশে) $\mSat{M}{p}$ হলে
    $\mSat{M}{q}$। কিন্তু $w_1$-তে $p$ সত্য ও $q$ অসত্য হওয়ায়
    $\mSat/{M}{p \lif q}$।
  \end{enumerate}
\end{proof}

\begin{prob}
  নিচে এমন ভাষার একটি মডেল $\mModel{M}$ দেওয়া আছে, যেখানে একমাত্র
  !!{propositional variable}s হলো $p_1$, $p_2$, $p_3$:
  \begin{center}
    \begin{tikzpicture}[modal]
      \node[world] (w1) [label={[align=right]left:\mTrue{p_1}\\\mFalse{p_2}\\\mFalse{p_3}}]
            {$w_1$} ;
      \node[world] (w2) [label={[align=right]right:\mTrue{p_1}\\\mTrue{p_2}\\\mFalse{p_3}},
        below right=of w1]  {$w_2$};
      \node[world] (w3) [label={[align=right]right:\mTrue{p_1}\\\mTrue{p_2}\\\mTrue{p_3}},
        above right=of w2] {$w_3$};
      \draw[reflexive above] (w3) to (w3);
      \draw[->] (w1) to (w2);
      \draw[->] (w2) to (w3);
      \draw[->] (w1) to (w3);
    \end{tikzpicture}
  \end{center}
  নিচের !!{formula}s ও স্কিমাগুলি কি মডেল $\mModel{M}$-এ, অর্থাৎ
  $\mModel{M}$-এর প্রত্যেক জগতে, সত্য? ব্যাখ্যা করো। এখানে পরমাণু
  চলটি উল্লিখিত তিনটির যেকোনোটি বোঝায়।
  \begin{enumerate}
  \item $p\lif \Diamond p$ ($p$ পরমাণু হলে);
  \item $!A\lif \Diamond !A$ ($!A$ যেকোনো সূত্র হলে);
  \item $\Box p \lif p$ ($p$ পরমাণু হলে);
  \item $\lnot p \lif \Diamond \Box p$ ($p$ পরমাণু হলে);
  \item $\Diamond \Box !A$ ($!A$ যেকোনো সূত্র হলে);
  \item $\Box \Diamond p$ ($p$ পরমাণু হলে)।
  \end{enumerate}
\end{prob}

\end{document}
